\chapter{凸分析与数学模型}\label{chap:convex-analysis}
\section{非线性规划的数学模型}\label{sec:modeling}

\begin{notebox}[引言]
非线性规划（Nonlinear Programming, NLP）比线性规划困难得多，主要原因在于：
\begin{itemize}[itemsep=0.2em]
    \item 没有统一的数学模型
    \item 每种方法有其特定的适用范围，不存在适用于各种NLP问题的一般算法
\end{itemize}
本节通过三个经典问题引入NLP的数学建模思想。
\end{notebox}

\subsection{选址问题}

\begin{conceptbox}[问题描述]
设有 $n$ 个市场，第 $j$ 个市场的位置为 $(a_j, b_j)$，其对某种货物的需求量为 $q_j\ (j=1,2,\ldots,n)$。现计划建立 $m$ 个仓库，第 $i$ 个仓库的存储容量为 $c_i\ (i=1,2,\ldots,m)$。试确定仓库的位置，使各仓库对各市场的运输量与路程乘积之和为最小。
\end{conceptbox}

\textbf{决策变量}：
\begin{itemize}
    \item 第 $i$ 个仓库的位置：$(x_i, y_i)\ (i=1,2,\ldots,m)$
    \item 第 $i$ 个仓库到第 $j$ 个市场的货物供应量：$w_{ij}$
\end{itemize}

\textbf{目标函数}（总运输成本最小）：
\[
\min\quad f = \sum_{i=1}^{m} \sum_{j=1}^{n} w_{ij} \sqrt{(x_i - a_j)^2 + (y_i - b_j)^2}
\]

\textbf{约束条件}：
\begin{enumerate}[label=(\arabic*)]
    \item 每个仓库的出货量不超过其容量：
    \[\sum_{j=1}^{n} w_{ij} \leq c_i,\quad i=1,2,\ldots,m\]
    \item 每个市场的需求被恰好满足：
    \[\sum_{i=1}^{m} w_{ij} = q_j,\quad j=1,2,\ldots,n\]
    \item 非负约束：
    \[w_{ij} \geq 0,\quad \forall i,j\]
\end{enumerate}

\begin{example}[选址问题建模]
\textbf{问题}：设有2个市场和1个仓库。市场1位于 $(0, 0)$，需求量 $q_1 = 5$；市场2位于 $(4, 0)$，需求量 $q_2 = 3$。仓库容量 $c_1 = 10$。求仓库最优位置。

\textbf{解}：设仓库位置为 $(x, y)$，运输量为 $w_1, w_2$。

数学模型为：
\begin{align*}
\min\quad & f = w_1 \sqrt{x^2 + y^2} + w_2 \sqrt{(x-4)^2 + y^2} \\
\text{s.t.}\quad & w_1 + w_2 \leq 10 \\
& w_1 = 5,\ w_2 = 3 \\
& w_1, w_2 \geq 0
\end{align*}

当 $w_1=5, w_2=3$ 固定时，最优点是两点之间的加权费马点。
\end{example}

\subsection{构件的表面积问题}

\begin{conceptbox}[问题描述]
一个半球形与圆柱形相接的构件，要求在构件体积一定的条件下，确定构件的尺寸使其表面积最小。
\end{conceptbox}

设圆柱形底半径为 $r$，高为 $h$。

\textbf{表面积}（目标函数）：
\[
S = 2\pi r^2 + 2\pi rh + \pi r^2 = 3\pi r^2 + 2\pi rh
\]
其中 $2\pi r^2$ 为半球面积，$2\pi rh$ 为圆柱侧面积，$\pi r^2$ 为圆柱底面积。

\textbf{体积约束}：
\[
V = \frac{2}{3}\pi r^3 + \pi r^2 h = V_0 \quad \text{(常数)}
\]

\textbf{数学模型}：
\[
\begin{aligned}
\min\quad & S = 3\pi r^2 + 2\pi rh \\
\text{s.t.}\quad & \frac{2}{3}\pi r^3 + \pi r^2 h = V_0 \\
& r \geq 0,\ h \geq 0
\end{aligned}
\]

\begin{example}[构件表面积问题求解]
\textbf{问题}：设体积 $V_0 = \frac{5\pi}{3}$，求使表面积最小的 $r$ 和 $h$。

\textbf{解}：由约束条件解出 $h$：
\[
h = \frac{V_0 - \frac{2}{3}\pi r^3}{\pi r^2} = \frac{5}{3r^2} - \frac{2r}{3}
\]

代入目标函数得无约束问题（关于 $r > 0$）：
\[
\min\quad S(r) = 3\pi r^2 + 2\pi r\left(\frac{5}{3r^2} - \frac{2r}{3}\right) = 3\pi r^2 + \frac{10\pi}{3r} - \frac{4\pi r^2}{3} = \frac{5\pi r^2}{3} + \frac{10\pi}{3r}
\]

求导：
\[
S'(r) = \frac{10\pi r}{3} - \frac{10\pi}{3r^2} = 0 \Rightarrow r^3 = 1 \Rightarrow r = 1
\]

验证二阶条件：
\[
S''(r) = \frac{10\pi}{3} + \frac{20\pi}{3r^3} > 0 \quad (r > 0)
\]

因此 $r^* = 1$ 是严格局部最小值点，代入得：
\[
h^* = \frac{5}{3} - \frac{2}{3} = 1,\quad S^* = \frac{5\pi}{3} + \frac{10\pi}{3} = 5\pi
\]
\end{example}

\subsection{木梁设计问题}

\begin{conceptbox}[问题描述]
把圆形木材加工成矩形横截面的木梁，要求木梁高度不超过 $H$，横截面的惯性矩不小于 $I$，且高度介于宽度与 $\alpha$ 倍宽度之间。如何确定木梁尺寸使成本最小？
\end{conceptbox}

设矩形横截面的高度为 $x_1$，宽度为 $x_2$，则圆形木材半径：
\[
R = \frac{1}{2}\sqrt{x_1^2 + x_2^2}
\]

\textbf{目标函数}（成本与横截面积成正比）：
\[
\min\quad f = \pi R^2 = \frac{\pi}{4}(x_1^2 + x_2^2)
\]

\textbf{约束条件}：
\begin{enumerate}[label=(\arabic*)]
    \item 高度上限：$x_1 \leq H$
    \item 惯性矩约束：$\frac{x_1^3 x_2}{12} \geq I$
    \item 高度宽度关系：$x_2 \leq x_1 \leq \alpha x_2$
    \item 非负约束：$x_1, x_2 \geq 0$
\end{enumerate}

\textbf{完整数学模型}：
\[
\begin{aligned}
\min\quad & f = \frac{\pi}{4}(x_1^2 + x_2^2) \\
\text{s.t.}\quad & x_1 \leq H \\
& \frac{x_1^3 x_2}{12} \geq I \\
& x_1 - x_2 \geq 0 \\
& x_1 - \alpha x_2 \leq 0 \\
& x_1, x_2 \geq 0
\end{aligned}
\]

\begin{notebox}[关键观察]
以上问题中，均含有决策变量的非线性函数，因此都属于\textbf{非线性规划问题}。
\end{notebox}

\subsection{Python代码：建模与可视化}

\begin{lstlisting}[language=Python, caption=选址问题可视化]
import numpy as np
import matplotlib.pyplot as plt
from scipy.optimize import minimize

# === 选址问题 ===
def facility_cost(vars, a, b, q):
    """
    vars = [x, y, w1, w2, ..., wn]
    (x,y): 仓库位置
    wi: 到第i个市场的运输量
    """
    x, y = vars[0], vars[1]
    w = vars[2:]
    cost = sum(wi * np.sqrt((x-ai)**2 + (y-bi)**2) 
               for wi, ai, bi in zip(w, a, b))
    return cost

# 示例：2个市场
a = np.array([0.0, 4.0])  # x坐标
b = np.array([0.0, 0.0])  # y坐标
q = np.array([5.0, 3.0])  # 需求量

# 可视化
x_range = np.linspace(-1, 5, 100)
y_range = np.linspace(-3, 3, 100)
X, Y = np.meshgrid(x_range, y_range)

# 固定w为需求量，只关于位置求值
Z = q[0]*np.sqrt(X**2 + Y**2) + q[1]*np.sqrt((X-4)**2 + Y**2)

plt.contour(X, Y, Z, levels=20)
plt.plot(a, b, 'ro', markersize=10, label='Markets')
plt.colorbar(label='Transport Cost')
plt.xlabel('x'); plt.ylabel('y')
plt.title('Facility Location Problem')
plt.legend(); plt.axis('equal')
plt.savefig('facility_location.png', dpi=150)
plt.show()
\end{lstlisting}

\begin{lstlisting}[language=Python, caption=构件表面积优化]
import numpy as np
from scipy.optimize import minimize_scalar

# === 构件表面积问题 ===
V0 = 5 * np.pi / 3

def surface_area(r):
    """S(r) = 5*pi*r^2/3 + 10*pi/(3*r)"""
    if r <= 0:
        return np.inf
    return 5*np.pi*r**2/3 + 10*np.pi/(3*r)

# 求解
result = minimize_scalar(surface_area, bounds=(0.01, 5))
r_opt = result.x
h_opt = 5/(3*r_opt**2) - 2*r_opt/3
S_opt = surface_area(r_opt)

print(f"Optimal r = {r_opt:.6f}")
print(f"Optimal h = {h_opt:.6f}")
print(f"Minimum S = {S_opt:.6f}")
print(f"Analytical: r*=1, h*=1, S*={5*np.pi:.6f}")
\end{lstlisting}

\section{非线性规划的一般数学模型}\label{sec:general-model}

\subsection{标准形式}

一般非线性规划的数学模型可表示为：
\[
\begin{aligned}
\min\quad & f(x) \\
\text{s.t.}\quad & g_i(x) \leq 0,\quad i=1,2,\ldots,m \\
& h_j(x) = 0,\quad j=1,2,\ldots,l
\end{aligned}
\]
其中 $x = (x_1, x_2, \ldots, x_n)^T \in \mathbb{R}^n$，$f: \mathbb{R}^n \to \mathbb{R}$，$g_i: \mathbb{R}^n \to \mathbb{R}$，$h_j: \mathbb{R}^n \to \mathbb{R}$。

\subsection{可行域与可行解}

\begin{definition}[可行域]
满足所有约束条件的解称为\textbf{可行解}（Feasible Solution）。所有可行解构成的集合称为\textbf{可行域}（Feasible Region）：
\[
\mathcal{X} = \{x \in \mathbb{R}^n \mid g_i(x) \leq 0,\ i=1,\ldots,m;\ h_j(x)=0,\ j=1,\ldots,l\}
\]
\end{definition}

\textbf{模型简记}：
\[
\min\ f(x),\quad x \in \mathcal{X}
\]

\subsection{问题分类}

\begin{conceptbox}[分类标准]
\begin{itemize}[itemsep=0.3em]
    \item \textbf{无约束问题}：$\mathcal{X} = \mathbb{R}^n$，即
    \[\min\ f(x),\quad x \in \mathbb{R}^n\]
    \item \textbf{有约束问题}：$\mathcal{X} \neq \mathbb{R}^n$
    \item \textbf{经典极值问题}：只有等式约束，可用拉格朗日乘子法化为无约束问题
\end{itemize}
\end{conceptbox}

\subsection{等价形式}

\textbf{不等式约束的不同写法}：
\begin{itemize}[itemsep=0.3em]
    \item $g_i(x) \leq 0$ 可改写为 $-g_i(x) \geq 0$
    \item 等式约束 $h_j(x) = 0$ 等价于两个不等式约束 $h_j(x) \leq 0$ 和 $-h_j(x) \leq 0$
\end{itemize}

\subsection{凸规划问题}

\begin{definition}[凸规划]
若目标函数 $f$ 和所有约束函数 $g_i$ 均为\textbf{凸函数}，等式约束 $h_j$ 为\textbf{仿射函数}，则称该问题为\textbf{凸规划问题}（Convex Programming）。
\end{definition}

\begin{property}[凸规划的优势]
凸规划问题的重要性质：
\begin{enumerate}[label=(\arabic*)]
    \item 局部最优解就是全局最优解
    \item 最优解集是凸集
    \item 若目标函数是严格凸函数且存在极小点，则全局最优解唯一
\end{enumerate}
\end{property}

\begin{example}[判断是否为凸规划]
\textbf{问题}：判断下列问题是否为凸规划：
\[
\begin{aligned}
\min\quad & f(x) = x_1^2 + 2x_2^2 \\
\text{s.t.}\quad & g_1(x) = x_1 + x_2 - 1 \leq 0 \\
& g_2(x) = -x_1 \leq 0
\end{aligned}
\]

\textbf{解}：
\begin{enumerate}[label=Step \arabic*:, leftmargin=2em]
    \item 检查 $f(x) = x_1^2 + 2x_2^2$：
    \[
    \nabla^2 f(x) = \begin{pmatrix} 2 & 0 \\ 0 & 4 \end{pmatrix}
    \]
    Hessian矩阵正定，故 $f$ 是严格凸函数。
    
    \item 检查 $g_1(x) = x_1 + x_2 - 1$：\textbf{仿射函数}（含常数项 $-1$），既是凸的也是凹的。
    
    \item 检查 $g_2(x) = -x_1$：线性函数（也是仿射），既是凸的也是凹的。
    
    \item 结论：所有约束为仿射（也是凸函数），目标函数严格凸，\textbf{这是凸规划问题}。
\end{enumerate}
\end{example}

\begin{example}[Logistic 损失函数的 Hessian 与凹性]\label{ex:logistic-hessian}
考虑损失函数
\[
l(\theta) = \sum_{i=1}^N \left[-\log(1 + \exp(-\theta^T u^{(i)})) - (1 - v^{(i)}) \theta^T u^{(i)}\right]
\]
其中 $u^{(i)}$ 为特征向量，$v^{(i)} \in \{0, 1\}$ 为标签。

记 $z_i = \theta^T u^{(i)}$，$\sigma_i = 1/(1 + e^{-z_i})$。梯度为
\[
\nabla l(\theta) = \sum_{i=1}^N (v^{(i)} - \sigma_i) u^{(i)}
\]
Hessian 矩阵为
\[
\nabla^2 l(\theta) = -\sum_{i=1}^N \sigma_i(1 - \sigma_i) u^{(i)} (u^{(i)})^T
\]
对任意向量 $a$：
\[
a^T \nabla^2 l(\theta) a = -\sum_{i=1}^N \sigma_i(1 - \sigma_i) (a^T u^{(i)})^2 \leq 0
\]
因此 $\nabla^2 l(\theta) \preceq 0$，$l(\theta)$ 是\textbf{凹函数}（即 $-l$ 是凸函数，可用于凸优化求解）。
\end{example}

\begin{lstlisting}[language=Python, caption=验证凸性（Python）]
import numpy as np

def check_convex_quadratic(H):
    """
    检查二次函数 f(x) = x^T H x 是否为凸函数
    通过验证Hessian矩阵是否半正定
    """
    eigenvalues = np.linalg.eigvals(H)
    print(f"Hessian矩阵: \n{H}")
    print(f"特征值: {eigenvalues}")
    if np.all(eigenvalues >= -1e-10):
        print("Hessian矩阵半正定 -> 函数是凸函数")
        if np.all(eigenvalues > 1e-10):
            print("Hessian矩阵正定 -> 函数是严格凸函数")
    else:
        print("Hessian矩阵不定 -> 函数不是凸函数")
    return eigenvalues

# 例题中的Hessian矩阵
H = np.array([[2, 0],
              [0, 4]])
eigenvals = check_convex_quadratic(H)
\end{lstlisting}

\section{非线性规划的直观理解}\label{sec:intuitive}

\subsection{等值线与图解法}

对于二元函数 $f(x_1, x_2)$，等值线定义为满足 $f(x_1, x_2) = c$（$c$ 为常数）的点的集合。

\begin{example}[图解法]
考虑非线性规划问题：
\[
\begin{aligned}
\min\quad & f(x_1, x_2) = (x_1 - 2)^2 + (x_2 - 1)^2 \\
\text{s.t.}\quad & x_1 + x_2 - 5 \leq 0 \\
& x_1, x_2 \geq 0
\end{aligned}
\]

\textbf{分析}：
\begin{itemize}
    \item 目标函数 $f$ 的等值线是以 $(2, 1)$ 为圆心的同心圆
    \item 无约束最优解为 $(2, 1)$，此时 $f(2,1) = 0$
    \item 检查 $(2, 1)$ 是否满足约束：$2 + 1 - 5 = -3 \leq 0$，满足！
    \item 因此无约束最优解就是约束最优解
\end{itemize}

若约束改为 $x_1 + x_2 - 2 \leq 0$，则 $(2, 1)$ 不可行。此时最优解在约束边界 $x_1 + x_2 = 2$ 与等值线相切处取得。
\end{example}

\begin{figure}[H]
\centering
\begin{tikzpicture}[scale=0.72]
    % 坐标轴
    \draw[->,thick] (-0.3,0) -- (5.5,0) node[right] {$x_1$};
    \draw[->,thick] (0,-0.3) -- (0,4.5) node[above] {$x_2$};
    % 等值线（同心圆，以(2,1)为圆心）
    \foreach \r/\c in {0.7/$1$, 1.4/$4$, 2.1/$9$, 2.8/$16$} {
        \draw[blue!50] (2,1) circle (\r);
        \node[blue!60,scale=0.7] at ({2+\r*0.707},{1+\r*0.707}) {\c};
    }
    % 约束线 x1+x2=5
    \draw[red,thick] (5,0) -- (0,5) node[pos=0.7,left,sloped,scale=0.75] {\small $x_1\!+\!x_2\!=\!5$};
    % 可行域
    \fill[green!12] (0,0) -- (5,0) -- (0,5) -- cycle;
    \draw[thick] (0,0) -- (5,0) -- (0,5) -- cycle;
    % 约束线 x1+x2=2
    \draw[red,dashed,thick] (2,0) -- (0,2) node[pos=0.5,left,sloped,scale=0.75] {\small $x_1\!+\!x_2\!=\!2$};
    % 最优点（无约束最优 = 约束最优）
    \fill[red] (2,1) circle (3pt) node[below right,scale=0.85] {\small $(2,1)$ 最优};
    % 切点（约束 x1+x2=2 时的最优）
    \fill[orange] (1.25,0.75) circle (2.5pt) node[below,scale=0.75] {\small 切点};
    % 梯度箭头
    \draw[->,thick,red!70!black] (2,1) -- (3.2,1.8) node[above right,scale=0.75] {$\nabla f$};
\end{tikzpicture}
\caption{非线性规划图解法：等值线（蓝色同心圆）与约束}
\label{fig:nlp-contour}
\end{figure}

\subsection{局部最优解与全局最优解}

\begin{definition}[局部最优解]
$x^*$ 是\textbf{局部最优点}，是指存在可行域中以 $x^*$ 为中心的邻域 $N(x^*, \delta)$，使得
\[
\forall x \in N(x^*, \delta) \cap \mathcal{X},\quad f(x^*) \leq f(x)
\]
如果不存在等号（对 $x \neq x^*$），则称为\textbf{严格局部最优解}。
\end{definition}

\begin{definition}[全局最优解]
若
\[
\forall x \in \mathcal{X},\quad f(x^*) \leq f(x)
\]
则 $x^*$ 为\textbf{全局最优解}。没有等号则为\textbf{严格全局最优解}。
\end{definition}

\begin{figure}[H]
\centering
\begin{tikzpicture}[scale=0.8]
    % 坐标轴
    \draw[->,thick] (-0.3,0) -- (7.5,0) node[right] {$x$};
    \draw[->,thick] (0,-0.3) -- (0,4.2) node[above] {$f(x)$};
    % 函数曲线（多峰函数）
    \draw[very thick,blue]
        (0.3,3.5) .. controls (0.8,2.5) and (1.2,1.8) .. (1.5,2)
        .. controls (1.8,2.2) and (2.2,0.5) .. (2.5,0.3)
        .. controls (3,0.2) and (3.5,1.5) .. (4,2.5)
        .. controls (4.5,3.2) and (5,2) .. (5.5,1)
        .. controls (6,0.5) and (6.5,1.5) .. (7,3);
    \node[blue,scale=0.8] at (7.2,3.2) {$f$};
    % 局部极小点
    \fill[orange] (1.5,2) circle (2.5pt) node[above,scale=0.7] {\small 局部极小};
    % 全局极小点
    \fill[red] (2.5,0.3) circle (3pt) node[below,scale=0.75] {\small \textbf{全局极小}};
    % 另一个局部极小
    \fill[orange] (5.5,1) circle (2.5pt) node[above,scale=0.7] {\small 局部极小};
    % 鞍点
    \fill[gray] (4,2.5) circle (2pt) node[above,scale=0.65] {\small 鞍点};
    % 局部极大
    \fill[green!50!black] (0.8,2.7) circle (2pt) node[above,scale=0.65] {\small 局部极大};
    % 邻域标注
    \draw[red,dashed] (2.5,0.3) circle (0.6);
    \node[red,scale=0.6,below] at (2.5,-0.3) {\small 邻域};
\end{tikzpicture}
\caption{局部最优解与全局最优解的对比（非凸函数可能有多个局部极小点）}
\label{fig:local-global}
\end{figure}

\begin{notebox}[LP vs NLP]
\begin{itemize}[itemsep=0.3em]
    \item \textbf{线性规划}：最优解必在可行域的极点达到，且一定是全局最优解
    \item \textbf{非线性规划}：最优解可能是可行域上的任一点，有局部和全局之分
    \item \textbf{一般NLP算法}往往求出的是局部最优解
    \item \textbf{凸规划}：极值点只有一个，既是局部最优点也是全局最优点
\end{itemize}
\end{notebox}

\subsection{一阶与二阶近似}

设目标函数 $f$ 在 $x^*$ 附近可微，考虑小扰动 $\Delta x$：

\textbf{一阶近似}（Taylor展开）：
\[
f(x^* + \Delta x) - f(x^*) \approx \nabla f(x^*)^T \Delta x
\]

\textbf{二阶近似}：
\[
f(x^* + \Delta x) - f(x^*) \approx \nabla f(x^*)^T \Delta x + \frac{1}{2} \Delta x^T \nabla^2 f(x^*) \Delta x
\]

\begin{conceptbox}[关键观察]
若 $x^*$ 是无约束局部极小点，则任意方向 $\Delta x$ 带来的一阶变化应非负：
\[
\nabla f(x^*)^T \Delta x \geq 0,\quad \forall \Delta x
\]
这要求 $\nabla f(x^*) = 0$（否则取 $\Delta x = -\nabla f(x^*)$ 可使左边为负）。
\end{conceptbox}

\begin{lstlisting}[language=Python, caption=等值线可视化]
import numpy as np
import matplotlib.pyplot as plt

# 定义函数及其梯度
def f(x, y):
    return (x - 2)**2 + (y - 1)**2

def grad_f(x, y):
    return np.array([2*(x - 2), 2*(y - 1)])

# 生成网格
x = np.linspace(-1, 4, 100)
y = np.linspace(-1, 4, 100)
X, Y = np.meshgrid(x, y)
Z = f(X, Y)

# 绘制等值线
plt.figure(figsize=(8, 6))
contours = plt.contour(X, Y, Z, levels=15, cmap='viridis')
plt.clabel(contours, inline=True, fontsize=8)

# 标记最优点
plt.plot(2, 1, 'r*', markersize=15, label='Optimal $(2,1)$')

# 绘制约束 x + y <= 2
x_con = np.linspace(-1, 4, 100)
y_con = 2 - x_con
plt.fill_between(x_con, -1, np.clip(y_con, -1, 4), 
                  alpha=0.2, color='red', label='$x_1+x_2 \\leq 2$')
plt.plot(x_con, y_con, 'r--', linewidth=2)

plt.xlabel('$x_1$'); plt.ylabel('$x_2$')
plt.title('Contour Plot and Feasible Region')
plt.legend(); plt.grid(True)
plt.axis('equal'); plt.xlim(-1, 4); plt.ylim(-1, 4)
plt.savefig('contour_plot.png', dpi=150)
plt.show()
\end{lstlisting}

\section{无约束问题的最优性条件}\label{sec:optimality-unconstrained}

\subsection{梯度与Hessian矩阵}

\begin{definition}[梯度]
多元函数 $f(x)$ 在 $x$ 处的\textbf{梯度}（Gradient）为：
\[
\nabla f(x) = \left(\frac{\partial f}{\partial x_1}, \frac{\partial f}{\partial x_2}, \ldots, \frac{\partial f}{\partial x_n}\right)^T
\]
\end{definition}

\textbf{梯度的几何意义}：
\begin{itemize}
    \item 梯度方向是函数值增加最快的方向
    \item 负梯度方向是函数值减小最快的方向（最速下降方向）
    \item 梯度垂直于等值面/等值线
\end{itemize}

\begin{definition}[Hessian矩阵]
多元函数 $f(x)$ 的二阶导数矩阵（Hessian矩阵）为：
\[
\nabla^2 f(x) = H(x) = \begin{pmatrix}
\frac{\partial^2 f}{\partial x_1^2} & \frac{\partial^2 f}{\partial x_1 \partial x_2} & \cdots & \frac{\partial^2 f}{\partial x_1 \partial x_n} \\[6pt]
\frac{\partial^2 f}{\partial x_2 \partial x_1} & \frac{\partial^2 f}{\partial x_2^2} & \cdots & \frac{\partial^2 f}{\partial x_2 \partial x_n} \\[6pt]
\vdots & \vdots & \ddots & \vdots \\[6pt]
\frac{\partial^2 f}{\partial x_n \partial x_1} & \frac{\partial^2 f}{\partial x_n \partial x_2} & \cdots & \frac{\partial^2 f}{\partial x_n^2}
\end{pmatrix}
\]
\end{definition}

\textbf{重要特例}：
\begin{itemize}
    \item 若 $f(x) = \frac{1}{2}x^T A x + b^T x + c$，$A$ 对称，则 $\nabla f(x) = Ax + b$，$\nabla^2 f(x) = A$
    \item 若 $f(x) = b^T x$，则 $\nabla f(x) = b$，$\nabla^2 f(x) = 0$
\end{itemize}

\begin{example}[计算梯度和Hessian矩阵]
\textbf{问题}：设 $f(x_1, x_2) = 3x_1^2 - 4x_1x_2 + 2x_2^2 + x_1 - 2x_2$，求梯度和Hessian矩阵。

\textbf{解}：
\begin{align*}
\frac{\partial f}{\partial x_1} &= 6x_1 - 4x_2 + 1 \\
\frac{\partial f}{\partial x_2} &= -4x_1 + 4x_2 - 2
\end{align*}

\textbf{梯度}：
\[
\nabla f(x) = \begin{pmatrix} 6x_1 - 4x_2 + 1 \\ -4x_1 + 4x_2 - 2 \end{pmatrix}
\]

\textbf{Hessian矩阵}：
\[
\nabla^2 f(x) = \begin{pmatrix} 6 & -4 \\ -4 & 4 \end{pmatrix}
\]

验证正定性：特征值为 $5 \pm \sqrt{17} > 0$，故正定，函数严格凸。
\end{example}

\subsection{Taylor展开}

\begin{property}[多元函数的Taylor展开]
\[
f(x + \Delta x) = f(x) + \nabla f(x)^T \Delta x + \frac{1}{2} \Delta x^T \nabla^2 f(x) \Delta x + O(\|\Delta x\|^3)
\]

\textbf{线性逼近}：$f(x + \Delta x) \approx f(x) + \nabla f(x)^T \Delta x$

\textbf{二次逼近}：$f(x + \Delta x) \approx f(x) + \nabla f(x)^T \Delta x + \frac{1}{2} \Delta x^T \nabla^2 f(x) \Delta x$
\end{property}

\subsection{一阶必要条件}

\begin{theorem}[一阶必要条件]\label{thm:first-order}
设多元函数 $f(x)$ 在 $x^*$ 点可微，若 $x^*$ 是 $f(x)$ 的局部极值点，则
\[
\nabla f(x^*) = 0
\]
满足 $\nabla f(x^*) = 0$ 的点称为\textbf{驻点}（Stationary Point）或\textbf{平稳点}。
\end{theorem}

\textbf{证明}（反证法）：
\begin{enumerate}[label=Step \arabic*:, leftmargin=2em]
    \item 假设 $\nabla f(x^*) \neq 0$
    \item 取方向 $d = -\nabla f(x^*)$，则对于充分小的 $\alpha > 0$：
    \[
    f(x^* + \alpha d) - f(x^*) \approx \alpha \nabla f(x^*)^T d = -\alpha \|\nabla f(x^*)\|^2 < 0
    \]
    \item 即 $f(x^* + \alpha d) < f(x^*)$，与 $x^*$ 是局部极小点矛盾
    \item 故 $\nabla f(x^*) = 0$
\end{enumerate}

\subsection{二阶必要条件}

\begin{theorem}[二阶必要条件]\label{thm:second-order-necessary}
设多元函数 $f(x)$ 在 $x^*$ 点二次可微，若 $x^*$ 是 $f(x)$ 的局部极小点，则
\[
\nabla f(x^*) = 0 \quad \text{且} \quad \nabla^2 f(x^*) \succeq 0 \ \text{（半正定）}
\]
\end{theorem}

\textbf{证明}（反证法）：
\begin{enumerate}[label=Step \arabic*:, leftmargin=2em]
    \item 由一阶必要条件，$\nabla f(x^*) = 0$ 已满足
    \item 假设 $\nabla^2 f(x^*)$ 不是半正定的，则存在负特征值
    \item 设 $d$ 为对应的特征向量，$\nabla^2 f(x^*) d = \lambda d$，$\lambda < 0$
    \item Taylor展开：
    \[
    f(x^* + \alpha d) - f(x^*) = \frac{\alpha^2}{2} d^T \nabla^2 f(x^*) d + o(\alpha^2) = \frac{\alpha^2 \lambda}{2} \|d\|^2 + o(\alpha^2) < 0
    \]
    对充分小的 $\alpha$，这与局部最优矛盾。
\end{enumerate}

\subsection{二阶充分条件}

\begin{theorem}[二阶充分条件 1]\label{thm:second-order-sufficient1}
设多元函数 $f(x)$ 在 $x^*$ 点二次可微，若
\[
\nabla f(x^*) = 0 \quad \text{且} \quad \nabla^2 f(x^*) \succ 0 \ \text{（正定）}
\]
则 $x^*$ 是 $f(x)$ 的\textbf{严格局部极小点}。
\end{theorem}

\textbf{证明}：由 $\nabla^2 f(x^*) \succ 0$，最小特征值 $\lambda_{\min} > 0$。对充分小的 $\|\Delta x\|$：
\[
f(x^* + \Delta x) - f(x^*) \approx \frac{1}{2} \Delta x^T \nabla^2 f(x^*) \Delta x \geq \frac{\lambda_{\min}}{2} \|\Delta x\|^2 > 0
\]

\begin{theorem}[二阶充分条件 2]\label{thm:second-order-sufficient2}
设多元函数 $f(x)$ 在 $x^*$ 的一个邻域内二次可微，若对该邻域内任意点：
\begin{itemize}
    \item $\nabla f(x^*) = 0$
    \item $\nabla^2 f(x)$ 半正定（不要求在 $x^*$ 处正定）
\end{itemize}
则 $x^*$ 是局部极小点。
\end{theorem}

\begin{notebox}[注意区别]
\begin{itemize}[itemsep=0.3em]
    \item 半正定性：特征值最小可为0
    \item 条件1要求 $x^*$ 处Hessian\textbf{正定}，结论更强（严格局部极小）
    \item 条件2允许Hessian在 $x^*$ 处为半正定，但在邻域内半正定也可保证局部极小
\end{itemize}
\end{notebox}

\subsection{利用最优性条件求解问题}

\begin{algobox}[求解步骤]
\begin{enumerate}[label=Step \arabic*:, leftmargin=2em]
    \item \textbf{求驻点}：解 $\nabla f(x) = 0$，得到所有候选点
    \item \textbf{二阶筛选}：计算各点的Hessian矩阵，滤除非半正定的点
    \item \textbf{判定严格极小}：若Hessian正定，确认为严格局部极小点
    \item \textbf{比较函数值}：检查所有满足必要条件的点，目标函数值最小的可能为全局极小点（需全局极小点存在的前提）
\end{enumerate}
\end{algobox}

\begin{example}[利用最优性条件求解]\label{ex:optimality1}
\textbf{问题}：利用最优性条件求解
\[
\min\ f(x, y) = x^4 + y^4 - 4xy
\]

\textbf{解}：

\textbf{Step 1}：求梯度并令其为零
\[
\nabla f = \begin{pmatrix} 4x^3 - 4y \\ 4y^3 - 4x \end{pmatrix} = \begin{pmatrix} 0 \\ 0 \end{pmatrix}
\]

即 $x^3 = y$，$y^3 = x$。代入得 $x^9 = x$，即 $x(x^8 - 1) = 0$。

四个驻点：
\[
x_1^* = (0, 0),\quad x_2^* = (1, 1),\quad x_3^* = (-1, -1),\quad x_4^* = (i, -i) \text{（舍去复数解）}
\]
实际驻点：$(0,0)$，$(1,1)$，$(-1,-1)$

\textbf{Step 2}：计算Hessian矩阵
\[
\nabla^2 f(x, y) = \begin{pmatrix} 12x^2 & -4 \\ -4 & 12y^2 \end{pmatrix}
\]

\textbf{Step 3}：逐点判断
\begin{itemize}
    \item 在 $(0, 0)$：$\nabla^2 f(0,0) = \begin{pmatrix} 0 & -4 \\ -4 & 0 \end{pmatrix}$，特征值为 $\pm 4$，不定，\textbf{不是极值点}（鞍点）
    \item 在 $(1, 1)$：$\nabla^2 f(1,1) = \begin{pmatrix} 12 & -4 \\ -4 & 12 \end{pmatrix}$，特征值为 $8, 16$，正定，\textbf{严格局部极小点}
    \item 在 $(-1, -1)$：$\nabla^2 f(-1,-1) = \begin{pmatrix} 12 & -4 \\ -4 & 12 \end{pmatrix}$，特征值为 $8, 16$，正定，\textbf{严格局部极小点}
\end{itemize}

\textbf{Step 4}：比较函数值
\[
f(1,1) = 1 + 1 - 4 = -2,\quad f(-1,-1) = 1 + 1 - 4 = -2
\]
两个点函数值相同，都是全局极小点。
\end{example}

\begin{example}[利用最优性条件求解（含半正定情形）]\label{ex:optimality2}
\textbf{问题}：利用最优性条件求解
\[
\min\ f(x, y) = x^2 - 2xy + y^2 = (x-y)^2
\]

\textbf{解}：

\textbf{Step 1}：求梯度
\[
\nabla f = \begin{pmatrix} 2x - 2y \\ -2x + 2y \end{pmatrix} = \begin{pmatrix} 0 \\ 0 \end{pmatrix}
\]
解得 $x = y$，即直线 $x = y$ 上所有点都是驻点。

\textbf{Step 2}：Hessian矩阵
\[
\nabla^2 f = \begin{pmatrix} 2 & -2 \\ -2 & 2 \end{pmatrix}
\]
特征值为 $4$ 和 $0$，半正定（非严格正定）。

\textbf{Step 3}：结论
由于 $f(x,y) = (x-y)^2 \geq 0$，且在直线 $x = y$ 上 $f = 0$，因此直线 $x = y$ 上所有点都是\textbf{全局最小点}。

注意：虽然Hessian不是严格正定的，但由于函数本身是非负的二次函数，仍能判定全局最优性。
\end{example}

\begin{lstlisting}[language=Python, caption=最优性条件验证]
import numpy as np
from sympy import *

# 定义符号
x, y = symbols('x y', real=True)

# 定义函数
f = x**4 + y**4 - 4*x*y

# 计算梯度
grad = [diff(f, x), diff(f, y)]
print("梯度:", grad)

# 求解驻点
critical_points = solve(grad, [x, y])
print("驻点:", critical_points)

# Hessian矩阵
H = Matrix([[diff(grad[0], x), diff(grad[0], y)],
            [diff(grad[1], x), diff(grad[1], y)]])
print("Hessian:", H)

# 逐点判断
for pt in critical_points:
    if pt[0].is_real and pt[1].is_real:
        H_val = H.subs([(x, pt[0]), (y, pt[1])])
        eigvals = H_val.eigenvals()
        print(f"点{pt}: Hessian={H_val}, 特征值={list(eigvals.keys())}")
\end{lstlisting}

\section{凸无约束问题的最优性条件}\label{sec:convex-optimality}

前面介绍的一般最优性条件只能用来研究局部最优解。对于一类特殊的函数——\textbf{凸函数}，可以获得全局最优性的充要条件。

\subsection{次梯度与次微分}

\begin{definition}[次梯度]
设 $f$ 是定义在非空开凸集 $C \subset \mathbb{R}^n$ 上的凸函数（不一定处处可微）。函数 $f$ 在 $x$ 处的\textbf{次梯度}（Subgradient）是一个向量 $g \in \mathbb{R}^n$，使得对 $C$ 内任意 $y$：
\[
f(y) \geq f(x) + g^T(y - x)
\]
所有次梯度的集合称为\textbf{次微分}（Subdifferential），记为 $\partial f(x)$。
\end{definition}

\begin{property}[次微分的性质]
\begin{itemize}[itemsep=0.3em]
    \item $\partial f(x)$ 是\textbf{非空}的、\textbf{凸的}、\textbf{紧的}
    \item $f$ 在 $x$ 可微，当且仅当 $\partial f(x) = \{\nabla f(x)\}$（单元素集合）
\end{itemize}
\end{property}

\begin{example}[次梯度计算]
\textbf{问题}：求 $f(x) = |x|$ 的次微分。

\textbf{解}：
\[
\partial f(x) = \begin{cases}
\{-1\}, & x < 0 \\[4pt]
[-1, 1], & x = 0 \\[4pt]
\{1\}, & x > 0
\end{cases}
\]
在 $x = 0$ 处，$f$ 不可微，次微分为区间 $[-1, 1]$。
\end{example}

\begin{example}[更多次微分计算]\label{ex:subdiff-more}
\textbf{(a)} $f(x_1, x_2, x_3) = \max\{|x_1|, |x_2|, |x_3|\} = \|x\|_\infty$ 在 $x = 0$ 处。

范数在原点的次微分等于其对偶范数单位球：
\[
\partial \|x\|_\infty \big|_{x=0} = \{g \in \mathbb{R}^3 : \|g\|_1 \leq 1\}
\]

\textbf{(b)} $f(x) = e^{|x|}$ 在 $x = 0$ 处（标量）。

令 $h(t) = e^t$，$h'(0) = 1$，且 $\partial |x|\big|_{x=0} = [-1, 1]$。由单调可微凸函数与凸函数复合的次微分链式法则：
\[
\partial f(0) = h'(0) \cdot \partial |x|\big|_{x=0} = [-1, 1]
\]

\textbf{(c)} $f(x_1, x_2) = \max\{x_1 + x_2 - 1,\ x_1 - x_2 + 1\}$ 在 $(1, 1)$ 处。

在 $(1, 1)$ 处两个仿射函数值均为 $1$，均为活跃函数。最大函数的次微分为活跃函数梯度的凸包：
\[
\partial f(1, 1) = \text{conv}\{(1, 1)^T, (1, -1)^T\} = \{(1, t)^T : t \in [-1, 1]\}
\]
\end{example}

\subsection{凸函数的一阶充要条件}

\begin{theorem}[凸函数的一阶充要条件]\label{thm:convex-first-order}
设 $f$ 定义在非空开凸集 $C \subset \mathbb{R}^n$ 上且可微，则 $f$ 为凸函数的充要条件是：
\[
f(y) \geq f(x) + \nabla f(x)^T(y - x),\quad \forall x, y \in C
\]
如果为严格不等式（$x \neq y$），则为严格凸函数。
\end{theorem}

\begin{figure}[H]
\centering
\begin{tikzpicture}[scale=0.85]
    % 坐标轴
    \draw[->,thick] (-0.3,0) -- (5.5,0) node[right] {$x$};
    \draw[->,thick] (0,-0.3) -- (0,4.2) node[above] {$f(x)$};
    % 凸函数曲线
    \draw[very thick,blue]
        (0.3,3.5) .. controls (1,1.5) and (2,0.5) .. (3,1) .. controls (3.5,1.5) and (4,2.5) .. (4.8,3.8);
    \node[blue,scale=0.8] at (5,3.9) {$f$};
    % 切点
    \fill[red] (2,0.5) circle (2.5pt) node[below left,scale=0.75] {\small $x$};
    % 切线
    \draw[red,thick,dashed] (0,1.5) -- (5,-1);
    \node[red,scale=0.75] at (4.5,-0.7) {\small 切线 $f(x)\!+\!\nabla f(x)^T(y\!-\!x)$};
    % 其他点 y
    \fill[black] (3.5,1.75) circle (2.5pt) node[above,scale=0.75] {\small $y$};
    \draw[dotted] (3.5,0) -- (3.5,1.75);
    \draw[dotted] (0,1.75) -- (3.5,1.75);
    % 切线上的对应点
    \fill[red!50] (3.5,-0.25) circle (2pt);
    \draw[<->,thick,green!50!black] (3.5,-0.25) -- (3.5,1.75);
    \node[green!50!black,scale=0.7,right] at (3.5,0.75) {\small $f(y)\!\geq$ 切线};
    % 凸组合标注
    \draw[<->,gray] (2,-0.6) -- (3.5,-0.6);
    \node[gray,scale=0.65,below] at (2.75,-0.6) {\small $y - x$};
\end{tikzpicture}
\caption{凸函数一阶充要条件：切线始终在函数图像下方}
\label{fig:convex-tangent}
\end{figure}

\textbf{等价表述}：$x^*$ 是凸函数 $f$ 的全局极小点，当且仅当
\[
0 \in \partial f(x^*) \quad \text{（即 $x^*$ 是驻点）}
\]

\begin{example}[复合优化问题的一阶条件]
\textbf{问题}：设 $f(x) = g(x) + \lambda \|x\|_1$，$g$ 光滑凸，$\lambda > 0$。求最优性条件。

\textbf{解}：

由于 $\|x\|_1$ 在 $x_i = 0$ 处不可微，使用次梯度：
\[
(\partial \|x\|_1)_i = \begin{cases}
\text{sign}(x_i), & x_i \neq 0 \\[-4pt]
[-1, 1], & x_i = 0
\end{cases}
\]

最优性条件：$0 \in \nabla g(x^*) + \lambda \partial \|x^*\|_1$

即对每个分量：
\[
(\nabla g(x^*))_i \in \begin{cases}
\{-\lambda\}, & x_i^* > 0 \\
\{\lambda\}, & x_i^* < 0 \\
[-\lambda, \lambda], & x_i^* = 0
\end{cases}
\]
这正说明了LASSO等稀疏优化中解的稀疏性来源：当 $|(\nabla g(0))_i| \leq \lambda$ 时，$x_i^* = 0$。
\end{example}

\subsection{凸函数的二阶充要条件}

\begin{theorem}[凸函数的二阶充要条件]\label{thm:convex-second-order}
设 $f$ 定义在非空开凸集 $C \subset \mathbb{R}^n$ 上且二阶可微，则 $f$ 为凸函数的充要条件是：
\[
\nabla^2 f(x) \succeq 0 \quad \text{（Hessian矩阵半正定，$\forall x \in C$）}
\]
若每点的Hessian矩阵正定，则 $f$ 为严格凸函数。
\end{theorem}

\begin{notebox}[重要区别]
\textbf{反之不成立}：严格凸函数只能导致Hessian半正定，不一定在每点都正定。例如 $f(x) = x^4$ 是严格凸的，但 $f''(0) = 0$。
\end{notebox}

\subsection{正定二次函数的性质}

对于正定二次函数：
\[
f(x) = \frac{1}{2}x^T A x + b^T x + c,\quad A \succ 0
\]

\begin{property}[二次函数的极值]
\begin{itemize}[itemsep=0.3em]
    \item \textbf{唯一极小点}：$x^* = -A^{-1}b$
    \item 等值面为以 $x^*$ 为中心的\textbf{同心椭球面族}
    \item 大多数优化方法从二次函数模型导出
\end{itemize}
\end{property}

\begin{definition}[二次终止性]
一个算法用于解正定二次函数的无约束极小时，有限步迭代可达最优解，则称该算法具有\textbf{二次终止性}。
\end{definition}

\textbf{二次终止性 = 共轭方向 + 精确一维搜索}

\subsection{凸函数的水平集}

\begin{definition}[水平集]
对于函数 $f$ 和实数 $\alpha$，水平集定义为：
\[
L_\alpha = \{x \in \dom(f) \mid f(x) \leq \alpha\}
\]
\end{definition}

\begin{property}[凸函数的水平集]
若 $f$ 是凸函数，则其任意水平集 $L_\alpha$ 都是凸集。反之不然！
\end{property}

\subsection{凸规划的全局最优性}

\begin{theorem}[凸规划的局部即全局]
设 $f$ 是定义在非空凸集 $X$ 上的凸函数，则 $f$ 的任意局部极小点就是其在 $X$ 上的\textbf{全局极小点}，且所有极小点形成一个凸集。
\end{theorem}

\textbf{证明}：
\begin{enumerate}[label=Step \arabic*:, leftmargin=2em]
    \item 假设 $x^*$ 为局部极小点，则存在邻域 $N(x^*, \delta)$ 使得
    \[f(x) \geq f(x^*),\quad \forall x \in N(x^*, \delta) \cap X\]
    \item 在 $X$ 中任取 $y$，连接 $x^*$ 和 $y$，存在 $\theta \in (0, 1)$ 使得
    \[x = \theta y + (1-\theta)x^* \in N(x^*, \delta) \cap X\]
    \item 由凸性：$f(x) \leq \theta f(y) + (1-\theta)f(x^*)$
    \item 由于 $f(x) \geq f(x^*)$，代入得：
    \[f(x^*) \leq \theta f(y) + (1-\theta)f(x^*) \Rightarrow f(x^*) \leq f(y)\]
    \item 由 $y$ 的任意性，$x^*$ 为全局极小点
\end{enumerate}

\begin{example}[凸规划的验证与求解]
\textbf{问题}：验证并求解
\[
\min\ f(x, y) = x^2 + xy + y^2 - 3x - 3y
\]

\textbf{解}：

\textbf{Step 1}：验证凸性
\[
\nabla^2 f = \begin{pmatrix} 2 & 1 \\ 1 & 2 \end{pmatrix}
\]
特征值为 $3$ 和 $1$，均正，Hessian正定，故 $f$ 严格凸。

\textbf{Step 2}：求驻点
\[
\nabla f = \begin{pmatrix} 2x + y - 3 \\ x + 2y - 3 \end{pmatrix} = \begin{pmatrix} 0 \\ 0 \end{pmatrix}
\]
解得 $x = y = 1$。

\textbf{Step 3}：由严格凸性，$x^* = (1, 1)$ 是\textbf{唯一全局最小点}。
\end{example}

\begin{lstlisting}[language=Python, caption=凸性验证与全局最优求解]
import numpy as np
from scipy.optimize import minimize

# === 凸函数验证与求解 ===
def f(vars):
    x, y = vars
    return x**2 + x*y + y**2 - 3*x - 3*y

def grad_f(vars):
    x, y = vars
    return np.array([2*x + y - 3, x + 2*y - 3])

def hess_f(vars):
    return np.array([[2, 1],
                     [1, 2]])

# 验证Hessian正定性
H = hess_f(None)
eigenvals = np.linalg.eigvals(H)
print(f"Hessian特征值: {eigenvals}")
print(f"最小特征值: {min(eigenvals):.4f} > 0, 严格凸")

# 求解
result = minimize(f, [0, 0], jac=grad_f, method='BFGS')
print(f"\\n最优解: x = {result.x}")
print(f"最优值: f* = {result.fun:.6f}")

# 验证解析解 x=y=1
print(f"\\n解析验证: f(1,1) = {f([1,1]):.6f}")
\end{lstlisting}

\section{解非线性规划问题的基本思路}\label{sec:solution-approach}

实际问题直接应用最优性条件往往很困难：导数不存在、Hessian计算复杂、方程 $\nabla f(x) = 0$ 求解困难等。因此通常采用\textbf{数值迭代方法}。

\subsection{迭代方法的基本框架}

\begin{algobox}[最优化方法的基本结构]
给定初始点 $x_0 \in \mathbb{R}^n$，$k = 0$，按以下规则迭代：
\begin{enumerate}[label=\arabic*., leftmargin=2em]
    \item \textbf{确定搜索方向} $d_k$：在 $x_k$ 点构造目标函数的下降方向
    \item \textbf{确定步长因子} $\alpha_k > 0$：使目标函数值有某种意义的下降
    \item \textbf{更新迭代点}：$x_{k+1} = x_k + \alpha_k d_k$
    \item 若满足某种\textbf{终止条件}，则停止迭代，得到近似最优解；否则 $k \leftarrow k+1$，返回步骤1
\end{enumerate}
\end{algobox}

\textbf{关键}：不同的步长因子 $\alpha_k$ 和不同的搜索方向 $d_k$ 构成了不同的优化方法。

\subsection{下降方向}

\begin{definition}[下降方向]
方向 $d$ 是函数 $f$ 在 $x$ 处的\textbf{下降方向}，若存在 $\bar{\alpha} > 0$ 使得：
\[
f(x + \alpha d) < f(x),\quad \forall \alpha \in (0, \bar{\alpha})
\]
\end{definition}

\textbf{判定条件}：若 $\nabla f(x)^T d < 0$，则 $d$ 是下降方向。

\textbf{证明}：由Taylor展开
\[
f(x + \alpha d) - f(x) = \alpha \nabla f(x)^T d + o(\alpha)
\]
当 $\nabla f(x)^T d < 0$ 且 $\alpha$ 充分小时，$f(x + \alpha d) < f(x)$。

\subsection{步长因子的确定——一维搜索}

\textbf{最优步长}通过求解一维优化问题得到：
\[
\alpha_k = \argmin_{\alpha > 0} f(x_k + \alpha d_k)
\]

\begin{theorem}[最优步长的正交性]
设 $x_{k+1} = x_k + \alpha_k d_k$ 是由精确线搜索得到的迭代点，则
\[
\nabla f(x_{k+1})^T d_k = 0
\]
即新点处的梯度与搜索方向正交。
\end{theorem}

\textbf{证明}：令 $\phi(\alpha) = f(x_k + \alpha d_k)$，则 $\alpha_k$ 满足 $\phi'(\alpha_k) = 0$，即
\[
\nabla f(x_k + \alpha_k d_k)^T d_k = \nabla f(x_{k+1})^T d_k = 0
\]

\subsection{终止条件}

常用的终止准则（Termination Criterion）：

\begin{conceptbox}[TC1：梯度准则]
\[
\|\nabla f(x_k)\| < \varepsilon_1
\]
表明梯度向量趋近于0，迭代序列收敛到平稳点。
\end{conceptbox}

\begin{conceptbox}[TC2：相邻迭代点准则]
\[
\|x_{k+1} - x_k\| < \varepsilon_2,\quad \text{或}\quad \frac{\|x_{k+1} - x_k\|}{\|x_k\|} < \varepsilon_2
\]
\end{conceptbox}

\begin{conceptbox}[TC3：函数值变化准则]
\[
|f(x_{k+1}) - f(x_k)| < \varepsilon_3,\quad \text{或}\quad \frac{|f(x_{k+1}) - f(x_k)|}{|f(x_k)|} < \varepsilon_3
\]
\end{conceptbox}

\begin{notebox}[Himmeblau终止准则（TC4）]
当TC2和TC3矛盾时（一个大一个小），采用组合准则：
\begin{itemize}
    \item 若 $|f(x_{k+1})| > \varepsilon$ 且 $\frac{|f(x_{k+1}) - f(x_k)|}{|f(x_{k+1})|} \leq \varepsilon_2$ 且 $\frac{\|\nabla f(x_{k+1})\|^2}{\|\nabla f(x_k)\|^2} < \varepsilon_3$
\end{itemize}

对于有一阶导数信息且收敛不太快的算法（如共轭梯度法），推荐TC1 + TC4结合使用，一般可取 $\varepsilon_1 = 10^{-4}$，$\varepsilon_2 = \varepsilon_3 = 10^{-5}$。
\end{notebox}

\subsection{收敛速度}

设迭代序列 $\{x_k\}$ 依某范数收敛到 $x^*$：

\begin{definition}[Q-收敛速度]
若存在实数 $p \geq 1$ 和与 $k$ 无关的正数 $c > 0$，使得
\[
\lim_{k \to \infty} \frac{\|x_{k+1} - x^*\|}{\|x_k - x^*\|^p} = c
\]
则称算法产生的点列具有 $p$ 阶Q-收敛速度：
\begin{itemize}[itemsep=0.2em]
    \item $p = 1$，$c \in (0, 1)$：\textbf{Q-线性收敛}
    \item $p = 1$，$c = 0$ 或 $p \in (1, 2)$：\textbf{Q-超线性收敛}
    \item $p = 2$：\textbf{Q-二阶（平方）收敛}
\end{itemize}
\end{definition}

\begin{conceptbox}[直观感受]
\begin{itemize}[itemsep=0.3em]
    \item \textbf{线性收敛}：每步误差约缩小为原来的 $c$ 倍（$c < 1$）
    \item \textbf{超线性收敛}：比任何线性速率都快
    \item \textbf{二阶收敛}：每步有效数字约翻倍
    \item \textbf{次线性收敛}：误差缩减越来越慢，如 $\frac{1}{k}$，$\frac{1}{k^2}$
\end{itemize}
\end{conceptbox}

\subsection{共轭方向法}

\begin{definition}[共轭方向]
设 $A$ 为 $n \times n$ 对称正定矩阵，非零向量 $d^{(1)}, d^{(2)}, \ldots, d^{(m)} \in \mathbb{R}^n$ 满足：
\[
(d^{(i)})^T A d^{(j)} = 0,\quad \forall i \neq j
\]
则称这 $m$ 个向量为 $A$-\textbf{共轭}的。若 $A = I$（单位矩阵），则共轭等价于正交。
\end{definition}

\begin{property}[共轭向量的线性无关性]
关于正定矩阵 $A$ 两两共轭的非零向量组必定线性无关。
\end{property}

\textbf{证明}：设 $\sum_{i=1}^{m} \beta_i d^{(i)} = 0$，左乘 $(d^{(j)})^T A$ 得 $\beta_j (d^{(j)})^T A d^{(j)} = 0$。由于 $A \succ 0$ 且 $d^{(j)} \neq 0$，故 $\beta_j = 0$。

\begin{theorem}[共轭方向法基本定理]
对于正定二次函数 $f(x) = \frac{1}{2}x^T A x + b^T x + c$，共轭方向法（GCD）至多经 $n$ 步精确线性搜索终止。
\end{theorem}

\begin{example}[共轭方向法迭代]
\textbf{问题}：用共轭方向法求解
\[
\min\ f(x, y) = x^2 + 2y^2
\]
取初始点 $x_0 = (2, 1)$，搜索方向 $d_0 = (-1, 0)$，$d_1 = (0, -1)$（关于 $A = \begin{pmatrix} 2 & 0 \\ 0 & 4 \end{pmatrix}$ 验证共轭性）。

\textbf{解}：

\textbf{Step 1}：沿 $d_0$ 搜索
\[
\min_\alpha\ f(2-\alpha, 1) = (2-\alpha)^2 + 2(1)^2
\]
求导：$-2(2-\alpha) = 0 \Rightarrow \alpha_0 = 2$
\[
x_1 = x_0 + 2d_0 = (2, 1) + 2(-1, 0) = (0, 1)
\]

\textbf{Step 2}：沿 $d_1$ 搜索
\[
\min_\alpha\ f(0, 1-\alpha) = 0 + 2(1-\alpha)^2
\]
求导：$-4(1-\alpha) = 0 \Rightarrow \alpha_1 = 1$
\[
x_2 = x_1 + 1d_1 = (0, 1) + (0, -1) = (0, 0) = x^*
\]

恰好 $n=2$ 步收敛到最优解！
\end{example}

\subsection{最优化算法分类概述}

\begin{conceptbox}[使用导数的方法]
\begin{itemize}[itemsep=0.3em]
    \item \textbf{最速下降法}：$d_k = -\nabla f(x_k)$，线性收敛，可能有锯齿现象
    \item \textbf{Newton法}：$d_k = -\nabla^2 f(x_k)^{-1} \nabla f(x_k)$，二阶收敛，但需Hessian正定
    \item \textbf{共轭梯度法}：利用梯度构造共轭方向，超线性收敛
    \item \textbf{拟Newton法（变尺度法）}：用近似矩阵 $G_k$ 代替Hessian的逆
    \begin{itemize}
        \item DFP公式、BFGS公式
    \end{itemize}
\end{itemize}
\end{conceptbox}

\begin{conceptbox}[不使用导数的方法（直接法）]
\begin{itemize}[itemsep=0.3em]
    \item \textbf{变量轮换法}：沿坐标轴方向轮换搜索
    \item \textbf{单纯形法}（Nelder-Mead）：通过反射、扩张、收缩操作
    \item 适用性强但收敛较慢
\end{itemize}
\end{conceptbox}

\begin{lstlisting}[language=Python, caption=最速下降法实现]
import numpy as np

def steepest_descent(f, grad_f, x0, alpha=0.1, 
                     tol=1e-6, max_iter=1000):
    """
    固定步长的最速下降法
    """
    x = x0.copy()
    history = [x.copy()]
    
    for k in range(max_iter):
        g = grad_f(x)
        if np.linalg.norm(g) < tol:
            break
        x = x - alpha * g
        history.append(x.copy())
    
    return x, f(x), k+1, history

# 测试
f = lambda x: x[0]**2 + 2*x[1]**2
grad_f = lambda x: np.array([2*x[0], 4*x[1]])

x_opt, f_opt, iters, hist = steepest_descent(
    f, grad_f, np.array([2.0, 1.0]), alpha=0.1)
print(f"最优解: x = {x_opt}")
print(f"最优值: f* = {f_opt:.8f}")
print(f"迭代次数: {iters}")
\end{lstlisting}

\section{有约束问题的最优性条件简介}\label{sec:constrained-kkt}

有约束优化问题的最优性条件（包括Lagrange乘子法和KKT条件）将在第~\ref{chap:constrained-theory}~章第~\ref{sec:ch6-kkt}~节中系统介绍。本节仅给出基本思路，详细内容请参见该章。
\begin{conceptbox}[有约束问题的最优性条件概览]
\begin{enumerate}[label=\arabic*., leftmargin=2em]
    \item \textbf{等式约束}：Lagrange乘子法（见第~\ref{chap:constrained-theory}~章第~\ref{sec:ch6-kkt}~节）
    \item \textbf{不等式约束}：KKT条件（见第~\ref{chap:constrained-theory}~章第~\ref{sec:ch6-kkt}~节）
    \item \textbf{凸规划}：KKT条件 = 充要条件（见第~\ref{chap:constrained-theory}~章定理~\ref{thm:kkt-sufficient}~）\end{enumerate}
\end{conceptbox}

\begin{notebox}[Fenchel 共轭函数]
函数 $f$ 的\textbf{Fenchel 共轭函数}定义为
\[
f^*(y) = \sup_x \{y^T x - f(x)\}
\]

\textbf{求解方法}：若 $f$ 可微且严格凸，令 $\nabla_x [y^T x - f(x)] = 0$，即 $y = \nabla f(x)$，解出 $x = x(y)$ 后代回 $f^*(y) = y^T x(y) - f(x(y))$。

\textbf{常见例子}：
\begin{itemize}[itemsep=0.2em]
    \item $f(x) = \frac{1}{2}\|x\|_2^2$ $\Rightarrow$ $f^*(y) = \frac{1}{2}\|y\|_2^2$
    \item $f(x) = \lambda\|x\|_1$ $\Rightarrow$ $f^*(y) = \begin{cases} 0, & \|y\|_\infty \leq \lambda \\ +\infty, & \text{otherwise} \end{cases}$
\end{itemize}

\textbf{与对偶的联系}：问题 $\min_x f(x) + g(Ax)$ 的 Fenchel 对偶为 $\max_y -f^*(-A^T y) - g^*(y)$。共轭函数是构造拉格朗日对偶、Fenchel 对偶和近端算法的基本工具。
\end{notebox}

\section{本章小结}\label{sec:ch3-summary}

\subsection{本章知识框架}

\begin{conceptbox}[非线性规划的数学模型——核心内容]
\begin{enumerate}[label=\arabic*., leftmargin=2em]
    \item \textbf{建模}：三个经典问题（选址、表面积、木梁设计）
    \item \textbf{一般模型}：$\min\ f(x)$，s.t. $g_i(x) \leq 0$，$h_j(x) = 0$
    \item \textbf{直观理解}：等值线、局部与全局最优
    \item \textbf{无约束最优性条件}：
    \begin{itemize}
        \item 一阶必要条件：$\nabla f(x^*) = 0$
        \item 二阶必要条件：$\nabla f(x^*) = 0$，$\nabla^2 f(x^*) \succeq 0$
        \item 二阶充分条件：$\nabla f(x^*) = 0$，$\nabla^2 f(x^*) \succ 0$
    \end{itemize}
    \item \textbf{凸规划的最优性条件}：
    \begin{itemize}
        \item 凸函数的充要条件：$\nabla^2 f(x) \succeq 0$（全局）
        \item 局部极小 = 全局极小
    \end{itemize}
    \item \textbf{迭代方法}：搜索方向 + 步长 + 终止条件 + 收敛速度
    \item \textbf{KKT条件}：有约束问题的最优性条件
\end{enumerate}
\end{conceptbox}

\subsection{关键公式汇总}

\begin{table}[htbp]
\centering
\caption{最优性条件汇总}
\begin{tabular}{@{}lll@{}}
\toprule
\textbf{条件类型} & \textbf{内容} & \textbf{性质} \\
\midrule
一阶必要 & $\nabla f(x^*) = 0$ & 驻点 \\
二阶必要 & $\nabla f(x^*) = 0$，$\nabla^2 f(x^*) \succeq 0$ & 局部极小 \\
二阶充分1 & $\nabla f(x^*) = 0$，$\nabla^2 f(x^*) \succ 0$ & 严格局部极小 \\
凸一阶 & $0 \in \partial f(x^*)$ & 全局极小（充要） \\
凸二阶 & $\nabla^2 f(x) \succeq 0$，$\forall x$ & 凸函数（充要） \\
KKT & 平稳性+可行性+对偶性+互补性 & 凸规划充要 \\
\bottomrule
\end{tabular}
\end{table}

\subsection{Python综合代码：求解无约束优化}

\begin{lstlisting}[language=Python, caption=综合Python求解示例]
import numpy as np
from scipy.optimize import minimize
import matplotlib.pyplot as plt

# ========== 综合求解示例 ==========

def solve_unconstrained(f, grad_f, hess_f, x0, method='BFGS'):
    """
    统一求解无约束优化问题
    """
    methods = {
        'BFGS': {'method': 'BFGS', 'jac': grad_f},
        'Newton-CG': {'method': 'Newton-CG', 
                      'jac': grad_f, 'hess': hess_f},
        'trust-ncg': {'method': 'trust-ncg',
                      'jac': grad_f, 'hess': hess_f},
        'CG': {'method': 'CG', 'jac': grad_f},
    }
    
    if method not in methods:
        raise ValueError(f"Unknown method: {method}")
    
    result = minimize(f, x0, **methods[method])
    return result

# === 测试函数：Rosenbrock函数 ===
def rosenbrock(x):
    """f(x,y) = (1-x)^2 + 100*(y-x^2)^2"""
    return (1 - x[0])**2 + 100 * (x[1] - x[0]**2)**2

def rosen_grad(x):
    dx = -2*(1 - x[0]) - 400*x[0]*(x[1] - x[0]**2)
    dy = 200*(x[1] - x[0]**2)
    return np.array([dx, dy])

def rosen_hess(x):
    dxx = 2 + 1200*x[0]**2 - 400*x[1]
    dxy = -400*x[0]
    dyy = 200
    return np.array([[dxx, dxy], [dxy, dyy]])

# 求解
print("=" * 50)
print("Rosenbrock函数优化")
print("=" * 50)

x0 = np.array([-1.0, 2.0])
for method in ['BFGS', 'Newton-CG', 'CG']:
    result = solve_unconstrained(rosenbrock, rosen_grad, 
                                  rosen_hess, x0, method)
    print(f"\\n{method}:")
    print(f"  最优解: x = {result.x}")
    print(f"  最优值: f* = {result.fun:.10f}")
    print(f"  迭代次数: {result.nit}")
    print(f"  函数调用: {result.nfev}")
    print(f"  梯度调用: {result.njev}")

# === 可视化收敛过程 ===
def rosenbrock_with_history(x):
    """带历史记录的Rosenbrock函数"""
    rosenbrock_with_history.iter_count += 1
    val = rosenbrock(x)
    rosenbrock_with_history.history.append((x.copy(), val))
    return val

rosenbrock_with_history.iter_count = 0
rosenbrock_with_history.history = []

result = minimize(rosenbrock_with_history, x0, 
                  method='BFGS', jac=rosen_grad)

history = rosenbrock_with_history.history
xs = [h[0][0] for h in history]
ys = [h[0][1] for h in history]
fs = [h[1] for h in history]

print(f"\\n收敛路径长度: {len(history)}")
print(f"初始函数值: {fs[0]:.4f}")
print(f"最终函数值: {fs[-1]:.10f}")
\end{lstlisting}

\subsection{进一步阅读}

\begin{notebox}[后续章节关联]
\begin{itemize}[itemsep=0.3em]
    \item \textbf{第四章}：线搜索方法（确定步长的具体算法）
    \item \textbf{第五章}：无约束最优化方法（最速下降、Newton、共轭梯度、BFGS等）
    \item \textbf{第六章}：有约束最优化方法（罚函数、内点法、SQP等）
    \item \textbf{凸分析}：深入理解次梯度、共轭函数、强凸性等概念
\end{itemize}
\end{notebox}

\subsection{关键概念速查}

\begin{table}[htbp]
\centering
\caption{关键术语中英文对照}
\begin{tabular}{@{}lll@{}}
\toprule
\textbf{中文} & \textbf{英文} & \textbf{说明} \\
\midrule
梯度 & Gradient & 一阶导数向量 \\
Hessian矩阵 & Hessian Matrix & 二阶导数矩阵 \\
驻点 & Stationary Point & $\nabla f = 0$ 的点 \\
次梯度 & Subgradient & 凸函数的推广梯度 \\
次微分 & Subdifferential & 所有次梯度的集合 \\
可行域 & Feasible Region & 满足约束的点集 \\
起作用约束 & Active Constraint & 紧约束 \\
拉格朗日乘子 & Lagrange Multiplier & KKT条件中的对偶变量 \\
精确线搜索 & Exact Line Search & 最优步长搜索 \\
二次终止性 & Quadratic Termination & $n$步收敛二次函数 \\
共轭方向 & Conjugate Direction & $d_i^T A d_j = 0$ \\
Q-线性收敛 & Q-linear Convergence & $\|e_{k+1}\| \leq c\|e_k\|$ \\
Q-超线性收敛 & Q-superlinear & $\|e_{k+1}\|/\|e_k\| \to 0$ \\
Q-二阶收敛 & Q-quadratic & $\|e_{k+1}\| \leq c\|e_k\|^2$ \\
\bottomrule
\end{tabular}
\end{table}

\newpage

% ==================== 第四章 一维搜索 ====================
